\documentclass[10pt,a4paper]{amsart}
\usepackage{amsmath,amssymb,amsthm,graphicx}
% Retained preamble for the Stacks Project chapter snapshot in this directory.
% The chapter's own preamble.tex is a copy of the shared Stacks Project preamble; the block
% below is sliced MECHANICALLY, line for line, from it. Nothing was added, removed or
% redefined; the wrapper class, the xr-hyper/hyperref wiring and the externaldocument list
% are omitted because the assembler supplies the document class.
%
% lines taken from preamble.tex: 3, 5, 9, 11, 12, 156, 157, 158, ... (51 lines total)
%
% package coverage restored from the chapter's own preamble (2026-10-13): xy, xy-2cell

\usepackage[all]{xy}
\xyoption{2cell}
\UseAllTwocells

\newenvironment{reference}{\comment}{\endcomment}
\newenvironment{slogan}{\comment}{\endcomment}
\newenvironment{history}{\comment}{\endcomment}
\theoremstyle{plain}
\newtheorem{theorem}[subsection]{Theorem}
\newtheorem{proposition}[subsection]{Proposition}
\newtheorem{lemma}[subsection]{Lemma}
\theoremstyle{definition}
\newtheorem{definition}[subsection]{Definition}
\newtheorem{example}[subsection]{Example}
\newtheorem{exercise}[subsection]{Exercise}
\newtheorem{situation}[subsection]{Situation}
\theoremstyle{remark}
\newtheorem{remark}[subsection]{Remark}
\newtheorem{remarks}[subsection]{Remarks}
\numberwithin{equation}{subsection}
\def\lim{\mathop{\mathrm{lim}}\nolimits}
\def\colim{\mathop{\mathrm{colim}}\nolimits}
\def\Spec{\mathop{\mathrm{Spec}}}
\def\Hom{\mathop{\mathrm{Hom}}\nolimits}
\def\Ext{\mathop{\mathrm{Ext}}\nolimits}
\def\SheafHom{\mathop{\mathcal{H}\!\mathit{om}}\nolimits}
\def\SheafExt{\mathop{\mathcal{E}\!\mathit{xt}}\nolimits}
\def\Sch{\mathit{Sch}}
\def\Mor{\mathop{\mathrm{Mor}}\nolimits}
\def\Ob{\mathop{\mathrm{Ob}}\nolimits}
\def\Sh{\mathop{\mathit{Sh}}\nolimits}
\def\NL{\mathop{N\!L}\nolimits}
\def\CH{\mathop{\mathrm{CH}}\nolimits}
\def\proetale{{pro\text{-}\acute{e}tale}}
\def\etale{{\acute{e}tale}}
\def\QCoh{\mathit{QCoh}}
\def\Ker{\mathop{\mathrm{Ker}}}
\def\Im{\mathop{\mathrm{Im}}}
\def\Coker{\mathop{\mathrm{Coker}}}
\def\Coim{\mathop{\mathrm{Coim}}}
\DeclareMathSymbol{\boxtimes}{\mathbin}{AMSa}{"02}
\def\QCohstack{\mathcal{QC}\!\mathit{oh}}
\def\Cohstack{\mathcal{C}\!\mathit{oh}}
\def\Spacesstack{\mathcal{S}\!\mathit{paces}}
\def\Quotfunctor{\mathrm{Quot}}
\def\Hilbfunctor{\mathrm{Hilb}}
\def\Curvesstack{\mathcal{C}\!\mathit{urves}}
\def\Polarizedstack{\mathcal{P}\!\mathit{olarized}}
\def\Complexesstack{\mathcal{C}\!\mathit{omplexes}}
\def\Pic{\mathop{\mathrm{Pic}}\nolimits}
\def\Picardstack{\mathcal{P}\!\mathit{ic}}
\def\Picardfunctor{\mathrm{Pic}}
\def\Deformationcategory{\mathcal{D}\!\mathit{ef}}
% lane-s11 extra preamble for sources whose statements carry \begin{reference} blocks.
% The blocks are purely bibliographic metadata (not mathematical content); the lane
% records the reference in the card's editorial notes, so the body skips them entirely.
\usepackage{comment}
\excludecomment{reference}
\usepackage{cleveref}
\crefname{originalreference}{original reference}{original references}
\Crefname{originalreference}{Original reference}{Original references}
\makeatletter
\newlabel{morphisms-example-scheme-theoretic-image}{{056A}{source}{}{sourcepointer-1}{}}
\newlabel{morphisms-example-scheme-theoretic-image@cref}{{[originalreference][0][]056A}{[1][0][]source}{}{}{}}
\makeatother
\title{Separating a closed subscheme from a fibre after a finite-presentation base change}
\author{The Stacks Project Authors}
\date{stacks-project:limits.tex:local-snapshot-20261012; GNU FDL 1.2 or later}
\begin{document}
\maketitle
\noindent\textit{Editorial note.} No omissions.\par\medskip
\noindent\textit{Allowed standard prerequisites (editorial).}\begin{itemize}
\item Scheme-theoretic image, affine computation. Morphisms, Example morphisms-example-scheme-theoretic-image
\end{itemize}\medskip
\ifdefined\hypertarget\hypertarget{sourcepointer-1}{}\fi
\noindent\textit{Source reference.} Original reference 056A: sources/stacks-chow/tags.txt, source lines 6703--6703. Cross-chapter standard prerequisite cited by the author inside the counted span; it is used as a pointer and its own proof is not part of this item.\par\medskip
\begin{lemma}
\label{lemma-separate}
Let $f : X \to S$ be a quasi-compact morphism of schemes.
Let $g : T \to S$ be a morphism of schemes.
Let $t \in T$ be a point and $Z \subset X_T$ be a closed
subscheme such that $Z \cap X_t = \emptyset$.
Then there exists an open neighbourhood
$V \subset T$ of $t$, a commutative diagram
$$
\xymatrix{
V \ar[d] \ar[r]_a & T' \ar[d]^b \\
T \ar[r]^g & S,
}
$$
and a closed subscheme $Z' \subset X_{T'}$ such that
\begin{enumerate}
\item the morphism $b : T' \to S$ is locally of finite presentation,
\item with $t' = a(t)$ we have $Z' \cap X_{t'} = \emptyset$, and
\item $Z \cap X_V$ maps into $Z'$ via the morphism $X_V \to X_{T'}$.
\end{enumerate}
Moreover, we may assume $V$ and $T'$ are affine.
\end{lemma}

\begin{proof}
Let $s = g(t)$. During the proof we may always replace $T$ by an
open neighbourhood of $t$. Hence we may also replace $S$ by an open
neighbourhood of $s$. Thus we may and do assume that $T$ and $S$ are affine.
Say $S = \Spec(A)$, $T = \Spec(B)$, $g$ is given by the
ring map $A \to B$, and $t$ correspond to the prime ideal
$\mathfrak q \subset B$.

\medskip\noindent
As $X \to S$ is quasi-compact and $S$ is affine we may write
$X = \bigcup_{i = 1, \ldots, n} U_i$ as a finite union of affine opens.
Write $U_i = \Spec(C_i)$. In particular we have
$X_T = \bigcup_{i = 1, \ldots, n} U_{i, T} =
\bigcup_{i = 1, \ldots n} \Spec(C_i \otimes_A B)$.
Let $I_i \subset C_i \otimes_A B$ be the ideal corresponding to the
closed subscheme $Z \cap U_{i, T}$. The condition that
$Z \cap X_t = \emptyset$ signifies that $I_i$ generates the
unit ideal in the ring
$$
C_i \otimes_A \kappa(\mathfrak q) =
(B \setminus \mathfrak q)^{-1}\left(
C_i \otimes_A B/\mathfrak q C_i \otimes_A B \right)
$$
Since $I_i (B \setminus \mathfrak q)^{-1}(C_i \otimes_A B) =
(B \setminus \mathfrak q)^{-1} I_i$ this means that $1 = x_i/g_i$
for some $x_i \in I_i$ and $g_i \in B$, $g_i \not \in \mathfrak q$.
Thus, clearing denominators we can find a relation of the form
$$
x_i + \sum\nolimits_j f_{i, j}c_{i, j} = g_i
$$
with $x_i \in I_i$, $f_{i, j} \in \mathfrak q$, $c_{i, j} \in C_i \otimes_A B$,
and $g_i \in B$, $g_i \not \in \mathfrak q$. After replacing $B$ by
$B_{g_1 \ldots g_n}$, i.e., after replacing $T$ by a smaller affine
neighbourhood of $t$, we may assume the equations read
$$
x_i + \sum\nolimits_j f_{i, j}c_{i, j} = 1
$$
with $x_i \in I_i$, $f_{i, j} \in \mathfrak q$, $c_{i, j} \in C_i \otimes_A B$.

\medskip\noindent
To finish the argument write $B$ as a colimit of finitely presented
$A$-algebras $B_\lambda$ over a directed set $\Lambda$.
For each $\lambda$ set
$\mathfrak q_\lambda = (B_\lambda \to B)^{-1}(\mathfrak q)$.
For sufficiently large $\lambda \in \Lambda$ we can find
\begin{enumerate}
\item an element
$x_{i, \lambda} \in C_i \otimes_A B_\lambda$ which maps to $x_i$,
\item elements $f_{i, j, \lambda} \in \mathfrak q_{i, \lambda}$
mapping to $f_{i, j}$, and
\item elements $c_{i, j, \lambda} \in C_i \otimes_A B_\lambda$
mapping to $c_{i, j}$.
\end{enumerate}
After increasing $\lambda$ a bit more the equation
$$
x_{i, \lambda} + \sum\nolimits_j f_{i, j, \lambda}c_{i, j, \lambda} = 1
$$
will hold. Fix such a $\lambda$ and set $T' = \Spec(B_\lambda)$.
Then $t' \in T'$ is the point corresponding to the prime $\mathfrak q_\lambda$.
Finally, let $Z' \subset X_{T'}$ be the scheme theoretic image of
$Z \to X_T \to X_{T'}$. As $X_T \to X_{T'}$ is affine, we can compute $Z'$
on the affine open pieces $U_{i, T'}$ as the closed subscheme associated
to $\Ker(C_i \otimes_A B_\lambda \to C_i \otimes_A B/I_i)$, see
Morphisms, Example \ref{morphisms-example-scheme-theoretic-image}.
Hence $x_{i, \lambda}$ is in the ideal defining $Z'$. Thus the last
displayed equation shows that $Z' \cap X_{t'}$ is empty.
\end{proof}

\end{document}
